\documentclass[UTF8,11pt,a4paper]{ctexart}
\usepackage[margin=2.45cm]{geometry}
\usepackage{amsmath,amssymb,amsthm,mathtools}
\usepackage{bm,booktabs,array,longtable}
\usepackage[colorlinks=true,linkcolor=blue,citecolor=blue,urlcolor=blue]{hyperref}
\usepackage{enumitem}

\newtheorem{theorem}{定理}[section]
\newtheorem{proposition}[theorem]{命题}
\newtheorem{lemma}[theorem]{引理}
\newtheorem{corollary}[theorem]{推论}
\theoremstyle{definition}
\newtheorem{definition}[theorem]{定义}
\newtheorem{problem}[theorem]{问题}
\theoremstyle{remark}
\newtheorem{remark}[theorem]{注}

\newcommand{\tr}{\operatorname{tr}}
\newcommand{\rank}{\operatorname{rank}}
\newcommand{\HS}{\mathrm{HS}}
\newcommand{\Id}{\mathrm{Id}}
\newcommand{\eps}{\varepsilon}
\newcommand{\cW}{\mathfrak W}
\newcommand{\spec}{\operatorname{spec}}
\newcommand{\Rea}{\operatorname{Re}}

\title{部分 Weil 配置、中心线零点比例与非零区域\\
\large 惯性、矩证书及数域实现的结构性框架}
\author{RH--Weil 结构研究项目}
\date{2026 年 9 月 1 日}

\begin{document}
\maketitle

\begin{abstract}
完整 Weil 正性迫使全部零点位于中心线，但对数域 zeta 函数而言，直接构造这种正性
几乎就是原问题本身。本文研究较弱而仍保留 Weil 几何的目标：有限压缩只需具有
函数方程给出的符号块、正确的一阶迹以及可控的二阶或高阶谱矩。我们定义
\emph{部分 Weil 配置}，从有限维 rank--trace 不等式推出稳健的抽象比例定理：
若归一化迹至少为 $\tau$，Hilbert--Schmidt 二阶矩至多为 $R$，则中心线上简单对象
的比例至少为 $4\tau-2-R$，互异对象的比例至少为 $(4\tau-R-1)/2$。

本文再给出两种升级机制。精确四矩 rank--trace--inertia 证书把 sine-kernel
矩数据对应的简单零点比例从 Christoffel 层级的 $13/18$ 提高到 $16/21$；
quadratic-channel Gram--Schur 分解则把所需四矩改写成背景、线性素数项和
二次素数项的有符号三通道账本。进一步，本文推导真实 Gabor 核的精确四循环
恒等式以及任意偶窗的配对对角泛函 $D_{22}(\psi)$；Montgomery--Taylor 窗给
$D_{22}=0.244589\ldots$，显示四矩目标窗不应由二矩优化单独决定。新的
signed-kernel 实现无条件消去 zero-side tail；纯素数 $2+2$ 又精确分裂为长度
$X$ 的 adjacent product Gram 与仍具 Farey $X^2$ 障碍的 alternating ratio Gram。
离线深度可见性把负谱迹转化为远离中心线的零点密度，并在负谱预算小于单个
离线对的可见深度时给出真正的非零区域。
这解释了为什么二阶矩足以证明正比例、却不能单独排除稀疏的离线零点。

Weil 显式公式和函数方程已经构造数域上的矩阵及符号块；当前素数侧二阶矩给出
Riemann zeta 与本原 Dirichlet $L$ 函数的 $0.67250\ldots$ 简单中心线比例。
进一步改进不必先证明完整四矩渐近：在平窗二矩下，只需把中心四矩一侧压到
$0.3275499\ldots$ 以下。本文把这一目标进一步分成 $2+2,3+1,4+0$ 素数共振族
和背景/taper/Gamma 余项的净预算。弱紧致性和凸分离可把“所有有限问题一致可解”
提升为非构造的全局响应，但不能凭空产生所需的一致算术预算。
\end{abstract}

\paragraph{结论状态。}
本文中的有限维线性代数定理均为无条件结论；应用到 zeta 的
$0.67250\ldots$ 数值是对 Alpöge--Furman 最新结果的结构化重述，而不是新的纪录。
高阶矩、深度可见性及 Brownian--Vaughan 响应与 zeta 的统一估计被标为待证输入。

\section{问题的放宽仍应保留 Weil 结构}

经典 Weil 判据把 Riemann 假设等价地写成一个 Hermitian 形式在全部测试函数上
半正定。有限域情形中，正极化、Hard Lefschetz 与函数方程共同把 Frobenius 谱压到
正确绝对值。数域中若直接假设同样的完整正性，所得“存在性定理”往往与 RH 循环。
因此更可操作的问题是：
\begin{quote}
有限压缩还保留哪些 Weil 型结构时，能够推出非零区域、中心线零点比例或互异性
比例？数域 zeta 的显式公式能否构造这些较弱配置，或者至少非构造地证明其存在？
\end{quote}

2026 年 Alpöge--Furman 的工作 \cite{AF26} 给出了决定性的例子：将 Weil
Hermitian 形式限制到约 $N(T,2T)$ 个调制测试函数上。中心线零点贡献半正定
rank-one 块；离线函数方程对贡献一个惯性为 $(1,1)$ 的块。素数侧则无条件给出
一阶迹与 Hilbert--Schmidt 二阶矩。rank--trace 不等式取代此前零点侧的逐项正性，
得到至少 $2/3$、经 Montgomery--Taylor 窗优化后至少 $0.67250\ldots$
的零点简单且在中心线上。

这说明，放宽目标并不要求回到与 Weil 路线无关的经典 mollifier 技术。零点比例
本身可由一个\emph{不完全正、但符号受函数方程控制}的 Weil 配置推出。本文抽离
该机制，并标出向更强结论升级时必需的新结构。

\section{部分 Weil 配置}

所有矩阵均为有限维复 Hermitian 矩阵。记 $n_+(A)$ 为 $A$ 的严格正特征值个数，
$n_-(A)$ 为严格负特征值个数，并令
\[
 \|A\|_{\HS}^2=\tr(A^2),\qquad
 A_-=\frac{|A|-A}{2},\qquad \mathcal{D}_-(A)=\tr A_-.
\]

\begin{definition}[有限部分 Weil 配置]
给定非负实数 $N,s,b,D,E_0,E_1$，一个
\[
 \cW=(G,P,Q;N,s,b,D;E_0,E_1)
\]
称为有限部分 Weil 配置，如果
\begin{align}
 &G=P+Q,\quad P\succeq0,\quad \rank P\le s,\quad \tr P\le s, \label{eq:pw1}\\
 &n_+(Q)\le b, \label{eq:pw2}\\
 &\tr P+2b\le N+E_0,\quad s+2b\le N+E_1, \label{eq:pw3}\\
 &D\ge s+b. \label{eq:pw4}
\end{align}
$N$ 是按重数计的总对象数；$s$ 是中心线上简单对象的容许上界；
$b$ 汇总中心线多重点和离线对的正惯性预算；$D$ 是互异对象数的下界账本。
$E_0,E_1$ 允许窗口边界、远尾或归一化误差。
\end{definition}

在 zeta 应用中，$P$ 只收集简单中心线零点。中心线多重点被放入 $Q$ 的正部分，
每个离线函数方程对给 $Q$ 至多一个正方向和一个负方向。因此
$b=s_2+p$，其中 $s_2$ 为中心线多重点数、$p$ 为离线对数；
$D\ge s_1+s_2+p$ 且 $N\ge s_1+2s_2+2p$。定义有意不引用 zeta，
因而同样适用于自对偶 $L$ 函数、相对行列式或任何带对称 divisor 的显式公式。

\section{rank--trace 引理与比例结构定理}

\begin{lemma}[rank--trace 不等式 {\cite[Lemma 3.2]{AF26}}]\label{lem:rt}
若 $P,Q$ 为 Hermitian 矩阵，$P\succeq0$、$\rank P\le r$ 且
$n_+(Q)\le b$，则
\begin{equation}\label{eq:rt}
 r\ge 2\tr P+4\tr Q-4b-\|P+Q\|_{\HS}^2.
\end{equation}
当 $P=\Pi_1$、$Q=2\Pi_2$，且 $\Pi_1,\Pi_2$ 是秩分别为 $r,b$
的正交投影时取等号。
\end{lemma}

\begin{proof}
写 $Q=Q_+-Q_-$，其中 $Q_\pm\succeq0$、$Q_+Q_-=0$ 且
$\rank Q_+\le b$。设 $p_i,n_i$ 分别为 $P,Q_-$ 的降序特征值。
von Neumann 迹不等式给出 $\tr(PQ_-)\le\sum_i p_i n_i$。
于是由 $x^2\ge2x-1$，
\[
 \|P\|_{\HS}^2-2\tr(PQ_-)+\|Q_-\|_{\HS}^2
 \ge \sum_i(p_i-n_i)^2
 \ge 2\tr P-r-4\tr Q_-.
\]
另一方面，对 $Q_+$ 的至多 $b$ 个正特征值使用 $x^2\ge4x-4$，得到
$\|Q_+\|_{\HS}^2\ge4\tr Q_+-4b$。再用 $\tr(PQ_+)\ge0$ 并相加，
即得 \eqref{eq:rt}。
\end{proof}

\begin{theorem}[部分 Weil 配置的比例定理]\label{thm:partial}
对任意有限部分 Weil 配置，
\begin{align}
 s&\ge4\tr G-\|G\|_{\HS}^2-2N-2E_0, \label{eq:sbound}\\
 D&\ge\frac{4\tr G-\|G\|_{\HS}^2-N-E_1}{2}. \label{eq:dbound}
\end{align}
右侧为负时只保留平凡下界 $0$。
\end{theorem}

\begin{proof}
在引理 \ref{lem:rt} 中取 $r=s$。由 $G=P+Q$，
\[
 s\ge4\tr G-2\tr P-4b-\|G\|_{\HS}^2.
\]
利用 $\tr P+2b\le N+E_0$ 即得 \eqref{eq:sbound}。
又因为 $\tr P\le s$，有 $3s+4b\ge4\tr G-\|G\|_{\HS}^2$。
而 $2D\ge2(s+b)=(3s+4b)-(s+2b)$，再用
$s+2b\le N+E_1$ 得 \eqref{eq:dbound}。
\end{proof}

\begin{corollary}[归一化缺陷账本]\label{cor:defect}
设一列配置满足 $N\to\infty$，
\[
 \tr G\ge\tau N,\qquad \|G\|_{\HS}^2\le RN,\qquad E_j\le\eps_jN.
\]
则
\begin{align}
 \frac{s}{N}&\ge\max\{0,\,4\tau-2-R-2\eps_0\}, \label{eq:norms}\\
 \frac{D}{N}&\ge\max\{0,\,\tfrac12(4\tau-R-1-\eps_1)\}. \label{eq:normd}
\end{align}
特别地，若 $\tau=1-\delta_{\rm tr}$、$R=1+\delta_2$ 且
$\eps_j=o(1)$，则
\[
 \frac{s}{N}\ge1-4\delta_{\rm tr}-\delta_2-o(1),\qquad
 \frac{D}{N}\ge1-2\delta_{\rm tr}-\frac{\delta_2}{2}-o(1).
\]
\end{corollary}

迹误差在简单中心线比例中被放大四倍，而二阶矩超额只支付一次。因此候选构造应
先确保临界密度框架的迹归一化，再优化 prime/Gamma/continuum 合并后的二阶矩。

\begin{corollary}[尖锐模型]\label{cor:sharp}
只使用 $\tr G$、$\|G\|_{\HS}^2$ 与块账本时，定理
\ref{thm:partial} 是尖锐的。令 $P=\Pi_1$、$Q=2\Pi_2$，
$\Pi_1\perp\Pi_2$，取 $\rank\Pi_1=2N/3$、$\rank\Pi_2=N/6$，则
\[
 \tr G=N,\quad \|G\|_{\HS}^2=\frac43N,\quad
 \frac{s}{N}=\frac23,\quad \frac{D}{N}=\frac56.
\]
\end{corollary}

所以要在相同信息层级上越过 $2-R$，不能仅改进线性代数常数；必须降低 $R$，
扩大测试函数 Fourier 支撑并获得新的 prime correlation 输入，或加入高阶谱信息。

\section{当前 zeta 实现与 0.67250\ldots}

Alpöge--Furman \cite{AF26} 的矩阵 $\widetilde G_T$ 是 Weil 形式在
$d\sim N(T,2T)$ 个临界密度 Gabor 测试函数上的压缩。其零点侧给出：
\begin{enumerate}[label=(Z\arabic*)]
\item 中心线简单零点产生半正 rank-one 项；
\item 离线函数方程对 $\{\rho,1-\bar\rho\}$ 的实矩阵项可写为
$2(aa^{\mathsf T}-bb^{\mathsf T})$，故正惯性至多为一；
\item 窗口远尾是迹范数 $o(1)$，且
$\tr\widetilde G_T=(1+o(1))N(T,2T)$。
\end{enumerate}
这些是 Weil 型部分：显式公式负责同一个 Hermitian form，函数方程负责符号块，
而不需要先知道零点在中心线上。

素数侧无条件计算
\[
 \|\widetilde G_T\|_{\HS}^2=(R(\psi)+o(1))N(T,2T).
\]
对指标窗 $R=4/3$，定理 \ref{thm:partial} 给 $2/3$ 和 $5/6$。
Montgomery--Taylor 窗给
\[
 R=c_{\rm MT}^{-1}
 =\frac12+\frac1{\sqrt2}\cot\frac1{\sqrt2},
\]
从而
\[
 2-R=0.67250\ldots,\qquad
 \frac{3-R}{2}=0.83625\ldots.
\]
该构造已无条件覆盖 Riemann zeta 和每个固定的本原 Dirichlet
$L(s,\chi)$ \cite{AF26}，且不用 mollifier、零密度估计或现成非零区域。

\begin{remark}
这项结果并没有说明剩余约 $32.75\%$ 的零点离开中心线；它们只是未被二阶矩证书
识别。$o(N)$ 个离线零点对迹和二阶矩可以不可见。因此比例定理本身不能通过取极限
自动变成 RH。
\end{remark}

\section{高阶矩与随机矩阵型谱证书}

二阶矩方法把全部谱压缩成两个数。若能从素数侧、随机矩阵模型或 Brownian 响应
得到更多混合矩，应直接逼近正半轴的谱指标函数，而不是把信息再次粗化为二阶矩。

\begin{theorem}[多项式正惯性证书]\label{thm:poly}
设 $\spec(G)\subset[-B,B]$。若实多项式
$p(x)=\sum_{k=0}^m a_kx^k$ 满足
\[
 p(x)\le0\quad(-B\le x\le0),\qquad
 p(x)\le1\quad(0<x\le B),
\]
则
\begin{equation}\label{eq:polyindex}
 n_+(G)\ge\tr p(G)=\sum_{k=0}^m a_k\tr(G^k).
\end{equation}
对部分 Weil 配置还有
\begin{align}
 s&\ge2\tr p(G)-N-E_1, \label{eq:polys}\\
 D&\ge\tr p(G). \label{eq:polyd}
\end{align}
\end{theorem}

\begin{proof}
对每个特征值 $\lambda$，有 $p(\lambda)\le\mathbf1_{\lambda>0}$，求和给
\eqref{eq:polyindex}。惯性次可加性给
$n_+(P+Q)\le\rank P+n_+(Q)\le s+b$。
于是 $D\ge s+b\ge n_+(G)$，且
$2n_+(G)\le2(s+b)\le N+E_1+s$。代入即可。
\end{proof}

\begin{corollary}[极限谱版本]
设经验谱测度 $\mu_T=d_T^{-1}\sum_j\delta_{\lambda_j(G_T)}$
弱收敛到支撑在 $[-B,B]$ 的概率测度 $\mu$，且所有 $G_T$ 的谱统一落在该区间。
则对每个满足定理 \ref{thm:poly} 条件的 $p$，
\[
 \liminf\frac{s_T}{N_T}
 \ge2\frac{d_T}{N_T}\int p\,d\mu-1-\limsup\frac{E_{1,T}}{N_T},
\quad
 \liminf\frac{D_T}{N_T}\ge\frac{d_T}{N_T}\int p\,d\mu.
\]
\end{corollary}

\subsection{四矩层级：13/18 与完整 rank--trace--inertia 改进}

Alpöge--Furman \cite[\S7.2(d)--(f)]{AF26} 指出：若归一化矩
\[
 m_k=N^{-1}\tr G^k\quad(0\le k\le4)
\]
取 sine-kernel Gram 极限值
\[
 (m_0,m_1,m_2,m_3,m_4)
 =\left(1,1,\frac43,2,\frac{13}{4}\right),
\]
则正惯性的 Christoffel 函数界
$n_+(G)/N\ge1-\Lambda_2(0)$、$\Lambda_2(0)=5/36$，
结合 $s/N\ge2n_+(G)/N-1$，给出
\[
 \frac{s}{N}\ge1-2\Lambda_2(0)=\frac{13}{18}=0.7222\ldots.
\]
这是\emph{四矩 Christoffel 层级}，不是随机矩阵路线的最终边界：
同一论文指出，若全部矩均满足相应假设，则该机制的比例上限为 $1$。

近期 working paper \cite{Notik26} 进一步保留 rank、trace 与 inertia 的完整账本，
而不先把它们粗化为单个正惯性，得到更强的四矩结论。下面把其核心有限维证书
纳入本框架；该线性代数部分可独立核验，但 zeta 的四矩算术假设仍是条件性的。

\begin{lemma}[谱删除账本]\label{lem:deletion}
设 $A$ 的特征值为 $\lambda_1\ge\cdots\ge\lambda_d$。存在
$P\succeq0$ 使
\[
 \rank P\le r,\qquad \tr P\le t,\qquad n_+(A-P)\le b
\]
当且仅当
\[
 \#\{i>b:\lambda_i>0\}\le r,\qquad
 \sum_{\substack{i>b\\\lambda_i>0}}\lambda_i\le t.
\]
\end{lemma}

\begin{proof}
必要性来自 Weyl 不等式
$\lambda_{i+b}(A)\le\lambda_i(P)+\lambda_{b+1}(A-P)\le\lambda_i(P)$，
再对正的左侧求和。反之，在 $A$ 的特征基中令
$P=\sum_{i>b,\,\lambda_i>0}\lambda_i v_iv_i^*$ 即可。
\end{proof}

引理说明归一化谱测度可分成三部分：至多质量 $b/N$ 的删除正谱
$\nu_d$；质量及一阶矩均至多 $s/N$ 的可覆盖正谱 $\nu_c$；
以及支撑在 $(-\infty,0]$ 的余项。

\begin{theorem}[精确四矩 rank--inertia 证书 {\cite{Notik26}}]
\label{thm:quartic}
假设 $m_0=m_1=1+o(1)$，令中心矩
\[
 b_2=m_2-1,\qquad b_4=m_4-4m_3+6m_2-3.
\]
若
\[
 1-2b_2+b_4>0,\qquad
 \sigma_*=\frac{b_2-b_4}{1-b_2}<\frac34,
\]
则部分 Weil 配置满足
\[
 \liminf\frac{s}{N}\ge
 R(b_2,b_4):=\frac{(1-b_2)^2}{1-2b_2+b_4},
\qquad
 \liminf\frac{D}{N}\ge\frac{1+R(b_2,b_4)}2.
\]
\end{theorem}

\begin{proof}[证书要点]
对 $\theta>0,\sigma<3/4$ 令
\[
 P_{\theta,\sigma}(x)
 =\theta(1-\sigma)^2+4\theta(1-\sigma)x
 -\theta\big((x-1)^2-\sigma\big)^2.
\]
写 $u=\theta(1-\sigma)$。直接求导和配方给
\[
 P_{\theta,\sigma}(x)\le
 \begin{cases}
 0,&x\le0,\\
 u(1-\sigma)+4ux,&x\in\mathbb R,
 \end{cases}
 \qquad
 P_{\theta,\sigma}(x)\le8u.
\]
将其分别积分于余项、$\nu_c$ 与 $\nu_d$，得到
\[
 \int P_{\theta,\sigma}\,d\nu
 \le u(5-\sigma)\frac{s}{N}+8u\frac{b}{N}.
\]
另一方面展开四次式，并用 $m_0=m_1=1$，
\[
 \int P_{\theta,\sigma}\,d\nu
 =u(5-\sigma)-\theta(\sigma^2-2b_2\sigma+b_4).
\]
再用 $s+2b\le N+o(N)$，可得
\[
 (1-\sigma)^2\left(1-\frac{s}{N}\right)
 \le\sigma^2-2b_2\sigma+b_4+o(1).
\]
在 $\sigma=\sigma_*$ 处优化即得第一个结论。对
$1-s/N-b/N$ 作同一估计并使用 $b/N\le1-s/N-b/N$，得到第二个结论。
\end{proof}

\begin{corollary}[同一四矩数据的更强条件比例]
若
$(m_1,m_2,m_3,m_4)=(1,\frac43,2,\frac{13}{4})+o(1)$，则
$b_2=1/3,b_4=1/4$，故
\[
 \frac{s}{N}\ge\frac{16}{21}=0.761904\ldots,\qquad
 \frac{D}{N}\ge\frac{37}{42}=0.880952\ldots.
\]
一个匹配的极端谱测度在
$1\pm1/(2\sqrt2)$ 各放质量 $8/21$，在 $2$ 和 $0$ 各放质量
$5/42$；它具有上述前四矩。因此在这些矩与账本内，$16/21$ 是尖锐值。
\end{corollary}

这使最值得优先尝试的算术目标发生变化。对平窗 $b_2=1/3$，若只想严格改善
当前无条件常数
\[
 \kappa=2-\left(\frac12+\frac1{\sqrt2}\cot\frac1{\sqrt2}\right),
\]
不必证明完整 Hardy--Littlewood 四矩渐近；只需一侧估计
\[
 b_4<\frac4{9\kappa}-\frac13
 =0.3275499074\ldots.
\]
预计值 $b_4=1/4$ 留有约 $31\%$ 的相对裕量。不过该 $b_4$ 是有效长度约
$T^2$、系数含 $(\Lambda*\Lambda)(m)m^{-1/2}$ 的正半定 Dirichlet 多项式均值；
普通 Montgomery--Vaughan 误差会比对角项大约多一个 $T$ 因子，仍需保留真实
算术相消。

随机矩阵型路线由三步组成：
\begin{enumerate}
\item 选择仍带零点侧函数方程块解释的 Weil 压缩 $G_T$；
\item 在素数侧证明 $\tr(G_T^k)$ 的极限或严格界，而不预设 RH；
\item 解单变量多项式 minorant 问题，将矩转为正惯性。
\end{enumerate}
随机矩阵启发可以建议极限谱和最优多项式，但真正的数论输入仍是第二步。通常的
GUE 零点间距猜想也不直接够用，因为这里需要的是压缩 Weil 矩阵的矩，而不是先
假定所有零点实后得到的 ordinate 点过程。

\begin{remark}[有理证书]
若没有统一算子范数界，可用有界有理函数替代多项式，例如由 resolvent
$(G-z)^{-1}$ 的矩阵元或 Stieltjes 变换构造
$r(x)\le\mathbf1_{x>0}$。这与本项目此前的 Cauchy/被动系统路线自然兼容，
也避免高次多项式在远谱处爆炸。
\end{remark}

\section{四矩的 quadratic-channel 因子化}

四矩算术困难不应被粗化成单个绝对值误差。下面的恒等式把它改写为三个
HS 通道的正半定 Gram，并保留背景与素数项的交叉相消。

\begin{theorem}[二次通道 Gram--Schur 分解]\label{thm:quadratic-channel}
设中心化压缩
\[
 H=G-I=A+V
\]
为有限维 Hermitian 算子，其中 $A$ 是不含素数振荡的背景，$V$ 是素数通道。
定义
\[
 C_0=A^2,\qquad C_1=AV+VA,\qquad C_2=V^2,
\quad
 K_{ij}=\tr(C_i^*C_j)\quad(0\le i,j\le2).
\]
则 $K\succeq0$，且
\begin{align}
 \tr H^4
 &=\|H^2\|_{\rm HS}^2
 =\|C_0+C_1+C_2\|_{\rm HS}^2 \notag\\
 &=K_{00}+2\Re K_{01}
   +(K_{11}+2\Re K_{02})+2\Re K_{12}+K_{22}.
 \label{eq:quadratic-channel-ledger}
\end{align}
令 $K_{<2}=(K_{ij})_{0\le i,j\le1}$、
$k=(K_{02},K_{12})^{\mathsf T}$，则
\[
 r_2:=K_{22}-k^*K_{<2}^{\dagger}k\ge0.
\]
若 $\mathcal P$ 是 Hilbert--Schmidt 空间到
$\operatorname{span}\{C_0,C_1\}$ 的正交投影，则
\[
 r_2=\|(I-\mathcal P)C_2\|_{\rm HS}^2,\qquad
 \tr H^4=\|C_0+C_1+\mathcal PC_2\|_{\rm HS}^2+r_2.
\]
\end{theorem}

\begin{proof}
$K$ 是三个向量 $C_0,C_1,C_2$ 的 Gram 矩阵，故为正半定。
又 $H^2=C_0+C_1+C_2$；展开其 Hilbert--Schmidt 范数即得
\eqref{eq:quadratic-channel-ledger}。Gram 矩阵的短接
$K_{22}-k^*K_{<2}^{\dagger}k$ 正是 $C_2$ 到前两个通道正交补的平方距离。
最后由勾股恒等式得到所述分解。
\end{proof}

式 \eqref{eq:quadratic-channel-ledger} 按 $V$ 的次数依次给出
\[
\begin{array}{c|ccccc}
\text{$V$-次数}&0&1&2&3&4\\ \hline
\text{贡献}&K_{00}&2\Re K_{01}&K_{11}+2\Re K_{02}
&2\Re K_{12}&K_{22}.
\end{array}
\]
因此不应单独上界 $\|V^2\|_{\rm HS}^2=K_{22}$；真正的 response-specific
目标是控制 Schur residual $r_2$ 及投影系数
$K_{<2}^{\dagger}k$。这与二矩路线中的 Vaughan--Brownian
公共通道相同：若先取各项绝对值，决定性的负交叉项会被人为删除。

\begin{lemma}[素数符号共振账本]\label{lem:prime-sign-ledger}
令
\[
 Q_T(t)=\sum_n c_n n^{it},\qquad V_T(t)=Q_T(t)+\overline{Q_T(t)},
\]
并令 $C_r(m)$ 为 $Q_T^r$ 的 Dirichlet 系数，$C_0(1)=1$。写
\[
 M_{r,s}=\frac1T\int_T^{2T}Q_T(t)^r\overline{Q_T(t)}^{\,s}\,dt.
\]
则
\[
 \frac1T\int_T^{2T}V_T(t)^4\,dt
 =6M_{2,2}+8\Re M_{3,1}+2\Re M_{4,0}.
\]
此外
\[
 M_{2,2}
 =\sum_{m,n}C_2(m)\overline{C_2(n)}
 K_T\!\left(\log\frac mn\right)\ge0,
\quad
 K_T(\omega)=\frac1T\int_T^{2T}e^{it\omega}\,dt.
\]
其对角项为 $\sum_m|C_2(m)|^2$；$2+2$ 非对角项以及
$3+1,4+0$ 两族均无固定符号。
\end{lemma}

\begin{proof}
展开 $(Q_T+\overline Q_T)^4$，并使用
$M_{s,r}=\overline{M_{r,s}}$。第二式是
$M_{2,2}=T^{-1}\int|Q_T^2|^2$ 的 Dirichlet 系数展开，非负性直接来自积分。
\end{proof}

对 zeta 的素数项，$c_n$ 的主尺度为 $\Lambda(n)n^{-1/2}$，所以
$C_2(m)$ 含 $(\Lambda*\Lambda)(m)m^{-1/2}$。这说明长度约 $T^2$
的困难不是一个任意双和，而是 quadratic coefficient map 的 Gram；
Vaughan 分解应作用于该系数映射，同时保留三通道的完整矩阵，而不是分别估计
$2+2,3+1,4+0$ 后相加绝对值。

\subsection{真实 Gabor 核的四循环与配对对角泛函}

二矩证明中的 Poisson--Gabor 恒等式可以无损地提升到四次迹，但把有限核替换为
平移不变核时，四矩需要新的边界估计。以下先把精确部分与尚待估计的部分分开。

\begin{theorem}[有限 Gabor 四循环恒等式]\label{thm:gabor-four-cycle}
沿用 Alpöge--Furman \cite[\S2]{AF26} 的记号，令
\[
 f_k(\tau)=\widehat\phi(\tau-\alpha_k),\qquad
 K(\tau,\tau')=\sum_{0\le k<d}f_k(\tau)f_k(\tau'),
\]
并令 $\mathcal B=\widetilde G+\widetilde E$ 为式 \cite[(2.11)]{AF26}
给出的全显式公式矩阵。则精确地
\[
 \tr\mathcal B^4
 =\frac1{a^4L^8}\int_{\mathbb R^4}
 \prod_{j=1}^4K(\tau_j,\tau_{j+1})
 \prod_{j=1}^4\nu_X(\tau_j)\,d\tau_1\cdots d\tau_4,
 \qquad \tau_5=\tau_1.
 \label{eq:finite-four-cycle}
\]
把每个 $K$ 形式替换为
$K_\infty(\tau,\tau')=L\Phi(\tau-\tau')$ 并把积分限制到
$I^4$，所得 bulk 泛函为
\[
 \mathcal C_4(\nu_X)
 =\frac1{a^4L^4}\int_{I^4}
 \prod_{j=1}^4\Phi(\tau_j-\tau_{j+1})
 \prod_{j=1}^4\nu_X(\tau_j)\,d\tau_1\cdots d\tau_4.
 \label{eq:bulk-four-cycle}
\]
\end{theorem}

\begin{proof}
将
$\mathcal B_{kk'}=(aL^2)^{-1}\int f_k(\tau)f_{k'}(\tau)\nu_X(\tau)d\tau$
代入
$\tr\mathcal B^4=\sum_{k_1,\ldots,k_4}
\mathcal B_{k_1k_2}\mathcal B_{k_2k_3}
\mathcal B_{k_3k_4}\mathcal B_{k_4k_1}$。
四个指标和分别组成四个 $K$，即得 \eqref{eq:finite-four-cycle}。
第二式只是代入 $K_\infty=L\Phi$ 后的正规化。
\end{proof}

\begin{proposition}[四循环的本征 Schatten--$4$ 稳定门]\label{prop:s4-gate}
在 $L^2(\mathbb R,d\tau)$ 上令 $S=M_{\operatorname{sgn}\nu_X}$，并令 $C_T\succeq0$
具有核
\[
 \frac{|\nu_X(\tau)|^{1/2}K(\tau,\tau')
 |\nu_X(\tau')|^{1/2}}{aL^2},
 \qquad J_T=C_T^{1/2}SC_T^{1/2}.
\]
把 $K$ 换成 $L\Phi$ 并乘
$\mathbf1_I(\tau)\mathbf1_I(\tau')$，同样定义
$C_{\infty,T},J_{\infty,T}$。则
\[
 \tr J_T^4=\tr\mathcal B^4,\qquad
 \tr J_{\infty,T}^4=\mathcal C_4(\nu_X).
\]
若
\[
 \mathcal C_4(\nu_X)=O(N),\qquad
 \|J_T-J_{\infty,T}\|_{\mathcal S_4}=o(N^{1/4}),
 \label{eq:s4-gate}
\]
则
\[
 \tr\mathcal B^4-\mathcal C_4(\nu_X)=o(N).
\]
\end{proposition}

\begin{proof}
对任意有限秩 $A,B$，迹的循环性与 Schatten Hölder 不等式给
\[
 |\tr A^4-\tr B^4|
 \le\|A-B\|_{\mathcal S_4}
 \bigl(\|A\|_{\mathcal S_4}^3
 +\|B\|_{\mathcal S_4}\|A\|_{\mathcal S_4}^2
 +\|B\|_{\mathcal S_4}^2\|A\|_{\mathcal S_4}
 +\|B\|_{\mathcal S_4}^3\bigr).
\]
令
$U_Tx=(\sqrt{a}\,L)^{-1}|\nu_X|^{1/2}\sum_kx_kf_k$，则
$\mathcal B=U_T^*SU_T$、$C_T=U_TU_T^*$。两次使用 $AB,BA$ 的非零谱
一致，得到 $\mathcal B,C_TS,J_T$ 的非零谱一致；bulk 同理。再取
$A=J_T,B=J_{\infty,T}$ 即得。
\end{proof}

命题 \ref{prop:s4-gate} 把 finite-to-bulk 问题定位为一个四次可和的核逼近，
并且只比较真正编码有符号谱的 Hermitian operators，而不要求更强的正包络差。

\begin{proposition}[zero-side tail 的四矩消去]\label{prop:zero-tail-fourth}
对固定 $\psi,\chi$，令 $m_L=1+\log L$。则
\[
 \|\widetilde G+\widetilde E\|
 \ll_{\psi,\chi}\frac{\sqrt Tm_L}{L},
\]
并且
\[
 \left|\tr(\widetilde G+\widetilde E)^4-\tr\widetilde G^4\right|
 \ll_{\psi,\chi}\frac{Tm_L^3}{L^3}=o(N).
 \label{eq:zero-tail-fourth}
\]
进一步，已有算子范数界下，一个可量化的 finite-to-bulk 充分条件是
\[
 \boxed{\quad
 \|J_T-J_{\infty,T}\|_{\rm HS}^2
 =o\!\left(\frac{L^3}{m_L^2}\right).
 \quad}
 \label{eq:boundary-hs-target}
\]
\end{proposition}

\begin{proof}
对 $g_x(\tau)=\sum_kx_k\widehat\phi(\tau-\alpha_k)$，Plancherel 与
$e^{2\pi iku/L}$ 在长度 $L$ 区间上的正交性给
$\int|g_x|^2\le2\pi L\|x\|^2$。在 $I$ 内用
$|\nu_X|\ll L+\sqrt X$；在 $I$ 外用 \cite[(5.3)--(5.4)]{AF26} 的
$\vartheta$ 包络及 $\int\vartheta^2=O(L)$、
$\int |r|\vartheta(r)^2dr=O(m_L)$，得到算子范数界。
再结合 \cite[Prop.~4.3]{AF26} 的
$\|\widetilde E\|_1\ll T^{-1/2}$ 和 trace-norm telescoping，得到
\eqref{eq:zero-tail-fourth}。

最后 $\|J_T-J_{\infty,T}\|=O(\sqrt Tm_L/L)$，而
$\|A\|_4^4\le\|A\|^2\|A\|_{\rm HS}^2$；式
\eqref{eq:boundary-hs-target} 因 $N\asymp TL$ 推出
$\|J_T-J_{\infty,T}\|_4^4=o(N)$，再应用命题 \ref{prop:s4-gate}。
\end{proof}

因此真正未闭合的边界只剩 $K_{\rm out}$ 与 $\tau$-domain 截断，而不是
zero-side tail。二矩的标量误差仍不能推出
\eqref{eq:boundary-hs-target}：取秩一 $A_T=(\sqrt Tm_L/L)P$、$B_T=0$，
其二矩差远小于现有保守误差尺度，但
$N^{-1}|\tr(A_T^4-B_T^4)|\asymp Tm_L^4/L^5\to\infty$。

这里还记录一个证据边界：按 \cite[\S5.2]{AF26} v2 当前展示的
$\mathcal E_2$ 末行直接正规化，只读出
\[
 O\!\left(\frac{N\log L}{L}\right),
\]
而命题 5.2 陈述的误差为
\[
 O\!\left(\frac{N\log L}{L^2}\right).
\]
这可能是展示证明省略了一步改进或排印差异；本文不据此断言原结论错误。
较弱界已足够其 $o(N)$ 二矩渐近，但不能把较强 rate 当作四矩 Schatten 输入。

若在算子层写 $\nu_X=A+P$，循环迹的十六个 words 精确合并为
\[
 \tr(A+P)^4
 =\tr A^4+4\tr(A^3P)+4\tr(A^2P^2)
 +2\tr(APAP)+4\tr(AP^3)+\tr P^4.
 \label{eq:cyclic-words}
\]
其中两个 $P$ 的六个 words 分成四个 adjacent words 与两个 alternating words；
这一区分在非交换 Gram 中不能抹去。

\begin{theorem}[纯素数 $2+2$ 的乘积/比值 Gram]\label{thm:product-ratio-gram}
定义
\[
 R_T(x)=\frac1{aL^2}\int_{\mathbb R}
 \mathbf f(\tau)\mathbf f(\tau)^*e^{ix\tau}\,d\tau,
 \qquad
 Z_T=-\frac1{2\pi}\sum_{n\le X}\frac{\Lambda(n)}{\sqrt n}R_T(\log n).
\]
则纯素数矩阵 $P_T=Z_T+Z_T^*$，且
\[
 [\tr P_T^4]_{2+2}
 =4\|Z_T^2\|_{\rm HS}^2+2\|Z_TZ_T^*\|_{\rm HS}^2.
 \label{eq:two-family-gram}
\]
若 $A=\|Z_T^2\|_{\rm HS}^2$、$B=\|Z_TZ_T^*\|_{\rm HS}^2$，则
\[
 B-A=\frac12\|[Z_T,Z_T^*]\|_{\rm HS}^2,
\]
故式 \eqref{eq:two-family-gram} 也等于
$6A+\|[Z_T,Z_T^*]\|_{\rm HS}^2$。

更细地，$Z_T^2$ 的 summands 可按乘积 $m=ab$ 聚成 PSD Gram，
$Z_TZ_T^*$ 的 summands 可按既约比值 $\rho=a/b$ 聚成 PSD Gram。
\end{theorem}

\begin{proof}
六个词在循环迹下分成四个 $ZZZ^*Z^*$ 和两个 $ZZ^*ZZ^*$；前者
等于 $\|Z^2\|_{\rm HS}^2$，后者等于 $\|ZZ^*\|_{\rm HS}^2$。
展开 $\|[Z,Z^*]\|_{\rm HS}^2$ 给 $2(B-A)$。最后分别写
\[
 Z_T^2=\sum_m\sum_{ab=m}b_ab_bR_aR_b,
\]
\[
 Z_TZ_T^*=\sum_{(p,q)=1}\sum_k
 b_{kp}\overline{b_{kq}}R_{kp}R_{kq}^*,
\]
即得两个 Gram 聚类。
\end{proof}

\begin{theorem}[偶窗口的配对 $2+2$ 对角泛函]
\label{thm:paired-window}
设 $\psi\ge0$ 为支撑在 $[-1/2,1/2]$ 的偶函数，以零延拓到全线，并令
$a=\int\psi>0$。对 $0\le r,s\le1$ 定义
\begin{align*}
 A_+(r,s)&=\int_{\mathbb R}\psi(u)^2\psi(u+r)\psi(u+s)\,du,\\
 A_-(r,s)&=\int_{\mathbb R}\psi(u)^2\psi(u+r)\psi(u-s)\,du,\\
 B(r,s)&=\int_{\mathbb R}\psi(u)\psi(u+r)\psi(u+s)\psi(u+r+s)\,du.
\end{align*}
则 bulk 泛函 \eqref{eq:bulk-four-cycle} 的纯素数 $P_X^4$ 项中，
由两两相等的乘法频率产生的对角贡献，除以 $N(T,2T)$ 后趋于
\[
 D_{22}(\psi)
 =\frac4{a^4}\int_0^1\!\!\int_0^1
 rs\bigl(A_+(r,s)+A_-(r,s)+B(r,s)\bigr)\,dr\,ds.
 \label{eq:paired-window-functional}
\]
更细地，两个 alternating sign words 的总贡献为
\[
 D_{\rm alt}(\psi)=\frac4{a^4}\iint rsA_+(r,s)\,dr\,ds,
\]
四个 adjacent sign words 的总贡献为
\[
 D_{\rm adj}(\psi)=\frac4{a^4}\iint
 rs\bigl(A_-(r,s)+B(r,s)\bigr)\,dr\,ds.
\]
\end{theorem}

\begin{proof}
写
$P_X(\tau)=-(2\pi)^{-1}\sum_{n\le X}a_n
\bigl(n^{i\tau}+n^{-i\tau}\bigr)$，$a_n=\Lambda(n)n^{-1/2}$。
四个符号中只有两个正号、两个负号能产生两两配对主项。对满足
$\lambda_1+\cdots+\lambda_4=0$ 的频率，令
$q=\phi^2$，对三个差变量作 Fourier 反演，得到
\[
 (2\pi)^3T\int_{\mathbb R}
 q(u)q(u+\lambda_1)q(u+\lambda_1+\lambda_2)
 q(u+\lambda_1+\lambda_2+\lambda_3)\,du
+\text{边界项}.
\]
再用 $q(u)=\psi(u/L)+o(1)$ 及弱极限
\[
 L^{-2}\sum_{n\le e^L}\frac{\Lambda(n)^2}{n}
 \delta_{\log n/L}\ \Longrightarrow\ r\,dr
\]
即可。三种 pair partitions 与四种 orientations 合并后分别产生
$A_+,A_-,B$。不同 pairings 的交集只含单素数幂退化项，其正规化为
$o(1)$。
\end{proof}

对指标窗，三个加权 overlap 积分依次为
$1/20,1/120,1/120$，所以
\[
 D_{\rm alt}(\psi_0)=\frac15,\qquad
 D_{\rm adj}(\psi_0)=\frac1{15},\qquad
 D_{22}(\psi_0)=\frac4{15}.
\]
这里 $1/5$ 与 $1/15$ 是两个 word families 的\emph{总和}，不是每个 word
各自的贡献。

对 Montgomery--Taylor 窗
$\psi_{\rm MT}(u)=\cos(\sqrt2u)\mathbf1_{|u|\le1/2}$，
\eqref{eq:paired-window-functional} 的收敛 Gauss 积分给
\[
 D_{22}(\psi_{\rm MT})\approx0.244589382034,
\quad
 D_{\rm alt}\approx0.178157229928,\quad
 D_{\rm adj}\approx0.066432152107.
\]
因此若同一窗口的真实中心四矩写成
$b_4=D_{22}(\psi_{\rm MT})+\mathcal R_{\rm net}$，
则在定理 \ref{thm:quartic} 的其他条件成立时，
\[
 \mathcal R_{\rm net}
 <b_2(\psi_{\rm MT})-D_{22}(\psi_{\rm MT})
 \approx0.082909914286
 \label{eq:mt-fourth-budget}
\]
已经足以严格改善该窗口自己的二矩比例。数值只评估显式窗口泛函；
它不是对 $\mathcal R_{\rm net}$ 的估计。

更一般地，记当前纪录为 $\kappa$。给定窗口的二阶中心矩 $b_2$ 时，
四矩证书超过 $\kappa$ 的一侧阈值是
\[
 b_4<B_\kappa(b_2)
 :=\frac{(1-b_2)^2}{\kappa}-1+2b_2,
 \label{eq:general-window-threshold}
\]
并仍须核对定理 \ref{thm:quartic} 的分母与 $\sigma_*$ 条件。
这说明窗口优化的正确目标不是只最小化二矩 $R(\psi)$，而是最大化
$B_\kappa(b_2(\psi))-D_{22}(\psi)$，同时控制未配对共振。
在余弦族 $\cos(cu)$ 中，有限扫描显示该预算向 $c\uparrow\pi$ 显著增大；
端点 $c=\pi$ 的形式值约为 $0.19827$，但这仍是选择下一候选窗的数值线索。

\begin{proposition}[adjacent 支撑缩减与端点窗边缘抑制]
\label{prop:path-support}
在 translation-invariant one-cut bulk 中，令
$r_n=\log n/L$。对 adjacent 次序 $+,+,-,-$，窗口 overlap 非零蕴含
\[
 r_a+r_b\le1,qquad ab\le X.
 \label{eq:adjacent-length}
\]
所以 adjacent product Gram 只具有长度 $X$；alternating 次序并无相应
乘积限制，其既约比值的 Farey 间距仍可达 $X^{-2}$。

更一般地，若四个 path shifts 的 span 为 $R\le1$，置 $\ell=1-R$，则
对 $\psi_c(u)=\cos(cu)\mathbf1_{|u|\le1/2}$、$0<c\le\pi$，
\[
 0\le\mathcal W_c
 \le v_c^2\ell+v_cc\ell^2+\frac{c^2}{6}\ell^3,
 \qquad v_c=\cos(c/2).
 \label{eq:cosine-edge}
\]
特别地，$\mathcal W_\pi\le\pi^2\ell^3/6$，而平窗只有
$\mathcal W_0=\ell$。
\end{proposition}

\begin{proof}
adjacent 部分和同时含 $0$ 与 $r_a+r_b$；长度一支撑要求全部部分和的
span 不超过一，即得 \eqref{eq:adjacent-length}。对第二式，把 overlap
交集写成 $u=-1/2+x$、$0\le x\le\ell$。达到两个极端 shift 的因子分别不超过
$v_c+cx$ 与 $v_c+c(\ell-x)$，另两个因子不超过一；积分即得
\eqref{eq:cosine-edge}。
\end{proof}

命题 \ref{prop:path-support} 把新的算术硬核精确缩到 alternating ratio Gram。
对任意固定 $\delta>0$，$m\le X^{1-\delta}$ 的 adjacent 部分已有普通
vector-valued large sieve 的 $O(T^{-\delta})$ 相对误差；只要再证明边缘
$X^{1-\delta}<m\le X$ 的 Gram mass 随 $\delta\downarrow0$ 消失，即可闭合
adjacent family。alternating family 则需要在真实 Vaughan response 子空间上
证明 Farey Gram 的一侧 Schur bound，不能使用 all-direction large sieve。

端点窗的 $O(\ell^3)$ 不代表逐点占优：所有 shifts 重合时，其正规化 overlap
为 $3\pi^4/128>1$。因此窗口筛选必须联合测量 bulk diagonal、边缘 path 和
alternating response norm。

\begin{proposition}[严格改善当前比例的一侧账本]\label{prop:fourth-ledger}
设 $b_2=1/3+o(1)$，并把真实 tapered 压缩的中心四矩写成
\[
 b_4=D_{22}+O_{22}+R_{31}+R_{40}+R_{\rm bg}+o(1),
\]
其中各符号分别表示 $2+2$ 对角、$2+2$ 非对角、$3+1$、
$4+0$ 与背景/taper/Gamma 余项。若存在 $\epsilon>0$ 和 $\delta_D$ 使
\[
 D_{22}\le\frac4{15}+\delta_D+o(1)
\]
及
\[
 O_{22}+R_{31}+R_{40}+R_{\rm bg}
 \le0.0608832407\ldots-\delta_D-\epsilon+o(1),
\]
则定理 \ref{thm:quartic} 给出严格大于当前
$0.6725007\ldots$ 的简单中心线零点比例。
\end{proposition}

\begin{proof}
两式给
$b_4<4/(9\kappa)-1/3$。在 $b_2=1/3$ 时，
$R(b_2,b_4)=4/[9(b_4+1/3)]>\kappa$，再应用定理
\ref{thm:quartic}。
\end{proof}

在平窗的形式 sine 账本中，两个 alternating words 合计 $1/5$，
四个 adjacent words 合计 $1/15$，故对角总贡献为 $4/15$；
预测的总四矩 $1/4$ 相当于净非对角修正 $-1/60$。
这些数只标定命题 \ref{prop:fourth-ledger} 的预算：
尚未把 alternating Farey Gram、$3+1,4+0$、Archimedean 背景、窗口边界
和截断尾统一压入该预算；但命题 \ref{prop:zero-tail-fourth} 已消去
zero-side tail，命题 \ref{prop:path-support} 又把 adjacent family 缩到长度 $X$。
有限无权 Dirichlet 多项式实验显示三族之间确有显著相消，但不构成 zeta
四矩估计。相应可执行账本见
\texttt{scripts/fourth\_moment\_channel\_audit.py}。

\section{Connes 非交换几何迹公式作为另一种配置载体}

Connes 的 1998 年论文 \cite{Connes98} 在 adele 类空间
$X=\mathbb A/k^*$ 上研究 idele 类群的作用，把临界零点解释为 absorption
spectrum，把可能的非临界零点解释为相对于中心线的 resonances。其有限地方集
$S$ 的迹公式得到证明；全局公式则通过 cutoff 投影 $Q_\Lambda$ 写成
\[
 \tr(Q_\Lambda U(h))
 =2h(1)\log'\Lambda+
 \sum_v\int_{k_v^*}'\frac{h(u^{-1})}{|1-u|}\,d^*u+o(1).
\]
在正特征全局域上，Connes 的 Theorem 5 证明此全局迹公式与所有
Grössencharakter $L$ 函数的 RH 等价；数域部分使用
Landau--Pollak--Slepian 的 prolate 投影几何处理 Archimedean cutoff。

这篇论文提供的是一个非常自然的\emph{非交换几何载体}，并不包含
$13/18$ 或其他零点比例常数。对本项目而言，正确的新问题不是再次假设其完整全局
迹公式，因为那会回到 RH 等价命题；而是建立一个较弱的“部分 Connes--Weil 迹
公式”：
\begin{enumerate}
\item 只对有限 Gabor/prolate family 控制 $\tr(Q_\Lambda U(h))$；
\item 证明函数方程共振在该压缩中具有所需惯性块；
\item 从局部或 $S$-local 几何估计前二矩、四矩的一侧上界或深度 frame 常数。
\end{enumerate}
若这三步可在不证明完整全局迹公式的情况下完成，Connes carrier 就可直接接入
定理 \ref{thm:partial}、\ref{thm:quartic} 或 \ref{thm:region}。这是一条与
显式公式 Gabor 压缩不同、但结构兼容的构造路线。

\section{从比例到非零区域：深度可见性}

离线块具有一个负方向，只是一个惯性陈述；它不保证有限测试族以统一强度看见
该方向。若零点贡献向量几乎落入压缩的核中，负特征值可以任意接近零。因此非零
区域需要比比例定理严格更强的输入。

\begin{definition}[深度可见性]\label{def:visibility}
记 $B_\eta$ 为给定高度窗内满足
$|\Rea\rho-\tfrac12|\ge\eta$ 的函数方程零点对数。
称 Hermitian 压缩 $G$ 对深离线零点具有 $(\eta,v,E_\eta)$-可见性，如果存在
子空间 $U_\eta$，满足
\[
 \dim U_\eta\ge B_\eta-E_\eta,\qquad
 \langle Gx,x\rangle\le-v(\eta)\|x\|^2
 \quad(x\in U_\eta),
\]
其中 $v(\eta)>0$。
\end{definition}

\begin{theorem}[负谱预算推出密度与非零区域]\label{thm:region}
若 $G$ 具有 $(\eta,v,E_\eta)$-可见性，则
\begin{equation}\label{eq:deepcount}
 B_\eta\le E_\eta+\frac{\mathcal{D}_-(G)}{v(\eta)}.
\end{equation}
因此：
\begin{enumerate}
\item 若 $E_\eta=o(N)$ 且
$\mathcal{D}_-(G)=o(v(\eta)N)$，则深离线零点对的相对密度趋于零；
\item 若 $E_\eta=0$ 且 $\mathcal{D}_-(G)<v(\eta)$，则 $B_\eta=0$；
\item 对变化的 $\eta=\eta(T)$，若第二条最终成立，则全部零点落在
$|\Rea s-\tfrac12|<\eta(T)$ 中。取
$\eta(T)=\tfrac12-c/\log T$ 时，这包含经典靠近 $\Rea s=1$
的非零区域表述。
\end{enumerate}
\end{theorem}

\begin{proof}
由 Courant--Fischer 极小极大原理，$G$ 至少有
$B_\eta-E_\eta$ 个特征值不大于 $-v(\eta)$。因此
$\mathcal{D}_-(G)\ge v(\eta)(B_\eta-E_\eta)$。
\end{proof}

\begin{proposition}[正负 synthesis 的可检验充分条件]\label{prop:synthesis}
设
\[
 G=CC^*-BB^*
\]
且存在 $r$ 维子空间 $U$ 与 $\alpha>\beta\ge0$，使得
\[
 \langle BB^*x,x\rangle\ge\alpha^2\|x\|^2,\qquad
 \langle CC^*x,x\rangle\le\beta^2\|x\|^2\quad(x\in U).
\]
则 $G|_U\preceq-(\alpha^2-\beta^2)\Id$。因而若
$r\ge B_\eta-E_\eta$，定义 \ref{def:visibility} 成立且
$v(\eta)=\alpha^2-\beta^2$。
\end{proposition}

命题 \ref{prop:synthesis} 把新输入分成两个量：负离线特征的下 frame bound
$\alpha$，以及正背景泄漏 $\beta$。最理想的情形是 $U\subset\ker C^*$，
此时只需控制负特征在商空间中的最小奇异值。更一般地可用 Schur/Feshbach 补
计算 $\alpha^2-\beta^2$。这正是 divisor-mode separation 的定量版本。

\begin{remark}[为什么二阶矩不够]
定理 \ref{thm:partial} 只使用 $n_+(Q)\le b$，而定理
\ref{thm:region} 要求负特征值远离零。有限个或 $o(N)$ 个极浅负方向可以同时满足
正确迹、正确二阶矩和任何固定正比例结论。因此从 $67.25\%$ 直接外推非零区域
在逻辑上不成立；必须证明可见性或等价的负谱分离。
\end{remark}

\section{构造性实现：数域上已经有什么、还缺什么}

\subsection{无条件可构造的三层}

对 Riemann zeta 与本原 Dirichlet $L$ 函数，以下结构已经可由经典数据构造：
\begin{enumerate}
\item \textbf{Hermitian 核：}Weil 显式公式把零点侧与
prime/Gamma/pole 侧写成同一个 sesquilinear form；
\item \textbf{符号块：}函数方程把离线零点配成
$\{\rho,1-\bar\rho\}$，其有限实压缩为 $(1,1)$ 型 pull-back；
\item \textbf{一、二阶矩：}临界密度测试族给正确迹；长度不超过 $T$ 的
Dirichlet 多项式均值给可计算的 Hilbert--Schmidt 矩。
\end{enumerate}
这三层已经足够实现定理 \ref{thm:partial}，并产生当前比例。

\subsection{三类真正开放的升级输入}

\begin{description}[leftmargin=3.7cm,style=nextline]
\item[更小的二阶矩常数]
在保持零点侧块账本和迹归一化的同时，使
$R=\|G\|_{\HS}^2/N$ 下降。带宽 $1$ 的现有窗类有自身天花板；
越过它需要支撑 $>1$ 的 prime pair-correlation 信息或新的预条件化压缩。

\item[高阶混合矩]
证明足够多的 $\tr G^k$ 或 resolvent 迹，使定理
\ref{thm:poly} 的证书优于二阶 rank--trace。随机矩阵模型在这里是设计工具，
而素数侧严格矩估计是证明输入。

\item[深度可见性与负谱预算]
证明深离线零点向量在测试族商空间中有统一 lower frame bound，并从算术侧控制
$\mathcal{D}_-(G)$。这是改善非零区域乃至最终排除离线零点所需的额外结构。
\end{description}

\subsection{响应预条件化与 Brownian--Vaughan 路线}

本项目此前把显式公式的算术 current 分成 Vaughan Type I、Type II 与连续通道，
并对 degree-two probe 构造了精确的 Brownian primitive response Gram。形式上，
若用可控响应 $R_T$ 对 Weil 压缩作预条件化，
\[
 G_T^{\rm corr}=G_T-\mathcal A_T R_T,
\]
理想目标是保持零点侧块账本与主迹，同时降低
$\|G_T^{\rm corr}\|_{\HS}^2$，或者在深离线商空间中增加负谱可见度。
这给出两个接口：
\begin{align}
 \Delta_2(R_T)
 &=\|G_T-\mathcal A_TR_T\|_{\HS}^2-\|G_T\|_{\HS}^2,\\
 v_T(\eta;R_T)
 &=\inf_{\substack{x\in U_{\eta,T}\\\|x\|=1}}
  -\langle(G_T-\mathcal A_TR_T)x,x\rangle.
\end{align}
前者直接改善比例，后者服务于非零区域。

已有笔记 195 的 exact channel Gram 表明 Type I、Type II 与 continuum 的交叉项
不能丢弃；对角能量可比完整能量大约高 $1.95$ 到 $3.36$ 倍。笔记 196 的
Cauchy translate 审计目前只捕获约 $2.1\%$ 到 $22.9\%$ 的响应能量。
这些数字只是有限诊断，尚未给出上述任一统一估计。特别是，Young/总变差界在自然
缩放下是 scale-neutral 的，不能单独使 $R$ 下降。

\section{非构造存在性：能做什么，不能做什么}

非构造方法的正确用途不是宣称某个完整 Weil 结构必存在；后者若不含独立的算术
约束，通常与 RH 等价。它们可以把无限个兼容约束降为有限可满足性，并把失败
转化为显式对偶证书。

\begin{theorem}[一致有限可满足性推出全局响应]\label{thm:compact}
设 $H$ 为 Hilbert 空间，$\{C_i\}_{i\in I}$ 是闭球
$\overline B_H(0,M)$ 中的弱闭集合。若任意有限 $J\subset I$ 都有
\[
 \bigcap_{i\in J}C_i\ne\varnothing,
\]
则 $\bigcap_{i\in I}C_i\ne\varnothing$。
特别地，若每个 $C_i$ 是范数闭凸集，则结论成立。
\end{theorem}

\begin{proof}
Hilbert 空间反射，故闭球在弱拓扑下紧。弱闭子集族具有有限交性质时，其总交非空。
闭凸集由 Hahn--Banach 分离定理知为弱闭。
\end{proof}

应用时，$H$ 可取 Brownian primitive response 空间，$C_i$ 编码有限组
prime/Gamma/continuum 通道的 affine 匹配、PSD block、Hilbert--Schmidt 球或
Schur complement 约束。定理 \ref{thm:compact} 说明：若所有有限通道问题都能
在同一个能量半径 $M$ 内完成，则存在一个同时完成全部约束的全局响应，即使没有
闭式公式。关键字是同一个 $M$；有限问题各自可解而最优范数趋于无穷时不能应用。

\begin{proposition}[凸锥的对偶存在性判据]\label{prop:dual}
设 $V$ 为有限维实 Hilbert 空间，$K\subset V$ 为闭凸锥。则
\[
 \operatorname{dist}(x,K)
 =\max_{\substack{y\in K^\circ\\\|y\|\le1}}\langle y,x\rangle,
 \qquad
 K^\circ=\{y:\langle y,k\rangle\le0\ \forall k\in K\}.
\]
因此 $\operatorname{dist}(x,K)\le\eps$ 当且仅当全部单位 polar separator
的响应不超过 $\eps$。
\end{proposition}

这允许把“存在一个正 Hodge 背景或低能响应”改写为对所有有限 rank-one packets
的统一不等式。它与本项目的一侧 atomic Hodge 定理兼容，也给计算机辅助搜索一个
清晰分工：数值 SDP 寻找最坏 separator，解析数论负责对该类 separator 给统一界。

\begin{remark}[非构造路线的边界]
紧致性和凸分离只传递已经证明的 uniform bounds，不制造它们；也不会自动给
深度可见性的 lower frame bound。若约束直接写成完整 Weil 形式半正，则其
有限可满足性本身可能已是 RH 强度。避免循环的标准是：约束必须只引用
prime/Gamma/continuum 数据及独立能量预算，而不能引用未知零点的位置。
\end{remark}

\section{可执行研究路线与停止条件}

\begin{longtable}{p{0.15\textwidth}p{0.25\textwidth}p{0.25\textwidth}p{0.17\textwidth}}
\toprule
路线 & 最小新引理 & 可改善的结论 & 停止或转向条件\\
\midrule
\endhead
二阶矩窗优化
& 在保迹、保块账本下证明 $R<R_{\rm MT}$
& 立即由 \eqref{eq:norms} 改善 $67.250\%$
& 证明目标窗类的对偶下界达到现常数\\
\addlinespace
支撑 $>1$
& 对更长 prime frequencies 的非对角均值给严格界
& 越过带宽一证书天花板
& 所需估计等价于未解 pair correlation 且无额外结构收益\\
\addlinespace
高阶矩或 RMT
& finite-to-bulk 四循环误差与未配对共振净预算
& 改善简单与互异比例
& 优化窗的 remainder budget 仍不能闭合\\
\addlinespace
深度可见性
& 商空间 lower frame bound $v_T(\eta)>0$
& 零密度、经典非零区域、最终可能排除
& $v_T(\eta)$ 因压缩核必然趋零且无可修复 frame\\
\addlinespace
Brownian 响应
& Type I/II/continuum 完整交叉 Gram 的 uniform Schur gain
& 降低 $R$ 或负谱预算
& gain 被尺度中性的能量下界完全抵消\\
\addlinespace
非构造紧致性
& 所有有限通道共同的能量半径
& 产生全局响应存在性
& 有限最优半径随通道数发散\\
\bottomrule
\end{longtable}

下一轮最值得推进的具体任务是：
\begin{enumerate}
\item 把 Alpöge--Furman 压缩矩阵写入本项目的
prime/continuum/Gamma response 记号，计算响应前后 $R$ 的精确 Schur 公式；
\item 在冻结有限模型上同时优化窗 $\psi$ 与 Brownian response，而不是只优化窗；
\item 证明 weighted boundary estimate
$\|J_T-J_{\infty,T}\|_{\rm HS}^2=o(L^3/(1+\log L)^2)$；
\item 证明 adjacent product Gram 的 edge tightness，并对 alternating Farey
Gram 构造 Vaughan response-specific Schur 证书；
\item 对离线测试零点 $\rho=\tfrac12+\eta+i\gamma$ 计算其 Gabor feature
在正背景商空间中的最小奇异值，建立 $v_T(\eta)$ 的模型尺度；
\item 将有限 response feasibility 写成 block PSD 与二阶锥问题，记录最优半径随
窗口和通道数的增长，以检验定理 \ref{thm:compact} 的 uniformity 前提。
\end{enumerate}

\section{结论}

部分 Weil 配置提供了一个位于完整 RH 与传统比例问题之间的中间结构：
\[
\boxed{\text{显式公式 Hermitian 核}
\ +\ \text{函数方程惯性块}
\ +\ \text{prime-side 谱矩}}
\]
一阶迹和二阶矩已经无条件给出 $67.250\ldots$；精确四矩证书与
quadratic-channel Gram 把比例改进压成 $2+2,3+1,4+0$ 共振族的一个净预算；
一般窗口的配对对角泛函又表明二矩最优窗未必是四矩最优窗。
进一步的 product/ratio 分解已把 $2+2$ 的普通长度通道与真正 Farey 硬核分开，
并把 finite-to-bulk 边界量化为一个 weighted Hilbert--Schmidt 目标。
负谱深度可见性则是非零区域所缺的结构。备选方向因而不是改做经典零点比例，
而是研究 Weil 配置的不完全正性如何量化地控制 divisor。

对数域存在性而言，显式公式已经构造配置的载体。真正要证明存在的不是任意
Hilbert 空间或形式极化，而是具有统一矩预算、响应能量和深度 frame 常数的算术
实现。弱紧致性可以完成最后的非构造拼接；最困难的部分仍是在
prime/Gamma/continuum 三者合并后证明这些 uniform constants。

\begin{thebibliography}{99}

\bibitem{AF26}
L. Alpöge and R. Furman,
\emph{More than two thirds of the zeros of the Riemann zeta function are
simple and on the critical line},
arXiv:2608.13637v2, 2026.
\url{https://arxiv.org/abs/2608.13637}.

\bibitem{Bombieri00}
E. Bombieri,
\emph{Remarks on Weil's quadratic functional in the theory of prime numbers I},
Rend. Mat. Acc. Lincei \textbf{11} (2000), 183--233.

\bibitem{Connes98}
A. Connes,
\emph{Trace formula in noncommutative geometry and the zeros of the
Riemann zeta function},
Selecta Math. (N.S.) \textbf{5} (1999), 29--106;
arXiv:math/9811068.

\bibitem{Notik26}
E. Notik,
\emph{A sharp four-moment inertia bound for the zeros of the Riemann
zeta function},
working paper, August 2026.
\url{https://www.academia.edu/171780663/}.

\bibitem{Conrey89}
J. B. Conrey,
\emph{More than two fifths of the zeros of the Riemann zeta function are on
the critical line},
J. Reine Angew. Math. \textbf{399} (1989), 1--26.

\bibitem{Levinson74}
N. Levinson,
\emph{More than one third of zeros of Riemann's zeta-function are on
$\sigma=1/2$},
Adv. Math. \textbf{13} (1974), 383--436.

\bibitem{Montgomery73}
H. L. Montgomery,
\emph{The pair correlation of zeros of the zeta function},
Proc. Sympos. Pure Math. \textbf{24} (1973), 181--193.

\bibitem{Weil52}
A. Weil,
\emph{Sur les formules explicites de la théorie des nombres premiers},
Comm. Sém. Math. Univ. Lund (1952), 252--265.

\end{thebibliography}

\end{document}
